\clearpage
% C057 | 示例修订：原三题保留题意，第三题换成更快的角度法；增加可行性反例。
\begin{examplebox}
  \textbf{\textcolor{CExample}{例1（基础：同支，直接倒数相加）}}

  \textbf{题目：} 双曲线 $\dfrac{x^2}{9}-\dfrac{y^2}{16}=1$ 的右焦点为 $F$，过 $F$ 的一条直线与\textbf{右支}交于不同两点 $A,B$，已知 $|AF|=3$，求 $|BF|$。

  \textbf{解：} $a=3,b^2=16$，同支公式给出
  \[
    \frac13+\frac1{|BF|}=\frac{2a}{b^2}=\frac38
    \quad\Longrightarrow\quad
    \frac1{|BF|}=\frac38-\frac13=\frac1{24}.
  \]
  \[
    \boxed{|BF|=24}.
  \]
  \textbf{核验：} $p=b^2/a=16/3$，$r_A=3>p/2=8/3$，另一段确实可以是正的有限长度。\textbf{识别信号：} “过焦点 + 同支 + 一段长”。

  \medskip
\end{examplebox}\clearpage\begin{examplebox}
  \textbf{\textcolor{CExample}{例2（变式：异支，先判断哪段更近）}}

  \textbf{题目：} 双曲线 $\dfrac{x^2}{4}-\dfrac{y^2}{16}=1$ 的右焦点为 $F$。经过 $F$ 的一条直线分别交右支于 $A$、左支于 $B$，已知 $|AF|=3$，求 $|BF|$。

  \textbf{解：} $a=2,b^2=16$，常数为 $2a/b^2=1/4$。对于\textbf{右焦点}，异支时右支交点 $A$ 离 $F$ 较近，因此\textbf{无需猜正负}：
  \[
    \frac1{|AF|}-\frac1{|BF|}=\frac14
    \quad\Longrightarrow\quad
    \frac1{|BF|}=\frac13-\frac14=\frac1{12}.
  \]
  \[
    \boxed{|BF|=12}.
  \]
  \textbf{几何核验：} $c=2\sqrt5$、$e=\sqrt5$。从焦半径得
  $x_A=\sqrt5,\ x_B=-2\sqrt5$；可取
  $A(\sqrt5,2)$、$B(-2\sqrt5,8)$，它们与 $F(2\sqrt5,0)$ 共线，且顺序为 $F-A-B$。

  \medskip
\end{examplebox}\clearpage\begin{examplebox}
  \textbf{\textcolor{CExample}{例3（综合：先求长度，再求斜率；精简原题）}}

  \textbf{题目：} 双曲线 $\dfrac{x^2}{9}-\dfrac{y^2}{16}=1$ 的右焦点为 $F$，过 $F$ 的一条直线与右支交于 $A,B$，且 $|AF|=2|BF|$。求直线斜率 $k$。

  \textbf{第一步（不用联立，求两段长）：} 设 $|BF|=t$，则 $|AF|=2t$。因两点同支，
  \[
    \frac1{2t}+\frac1t=\frac38
    \quad\Longrightarrow\quad t=4,
    \qquad\boxed{|AF|=8,\quad |BF|=4}.
  \]
  \textbf{第二步（最快：由方向余弦求斜率）：} $e=5/3$，$p=16/3$。取 $F$ 指向 $A$ 的方向与 $x$ 轴正向夹角为 $\theta$，由右支焦半径的倒数式
  \[
    \frac18=\frac{1-e\cos\theta}{p}
    \quad\Longrightarrow\quad
    e\cos\theta=1-\frac p8=\frac13,
    \quad \cos\theta=\frac15.
  \]
  故直线斜率的绝对值为
  \[
    |k|=|\tan\theta|
       =\frac{\sqrt{1-\cos^2\theta}}{|\cos\theta|}
       =\frac{\sqrt{24/25}}{1/5}=2\sqrt6.
  \]
  由于上下关于 $x$ 轴对称，
  \[
    \boxed{k=\pm2\sqrt6}.
  \]
  \textbf{第三步（常规坐标复核）：} $r=ex-a$ 给出
  $x_A=(8+3)/(5/3)=33/5$，$x_B=(4+3)/(5/3)=21/5$。
  代入双曲线可取
  \[
    A\left(\frac{33}{5},\frac{16\sqrt6}{5}\right),\qquad
    B\left(\frac{21}{5},-\frac{8\sqrt6}{5}\right).
  \]
  两点都在双曲线上且 $F=(5,0)$，计算
  \[
    \frac{y_A}{x_A-5}=2\sqrt6,
    \qquad \frac{y_B}{x_B-5}=2\sqrt6.
  \]
  故它们与 $F$ 共线；同时将纵坐标全部变号得到另一斜率。\textbf{经验：} 比起逐点解 $y$，求斜率时优先用\textbf{方向余弦}。

  \medskip
\end{examplebox}\clearpage\begin{examplebox}
\begin{center}\begin{tikzpicture}[font=\small,>=stealth]
\draw[rounded corners=3pt,gray!45] (0,0) rectangle (14.2,12.8);
\node at (7.1,12.2) {同支时：两块倒数都为正，才能拼成固定总量 $3/8$};\node at (7.1,10.75) {例1：同支 $AF=3$，另一段 $BF=24$};
\draw[gray,dashed] (2.2,9.4) rectangle (8.2,10.299999999999999);\fill[blue!30] (2.2,9.6) rectangle (7.533333333333333,10.1);
\fill[orange!45] (7.533333333333333,9.6) rectangle (8.2,10.1);
\draw[violet,dashed] (8.2,9.35)--(8.2,10.45);\node[anchor=west] at (10.7,9.85) {正的余量};
\node at (5.2,9.049999999999999) {$1/3+1/24=3/8$};
\node at (7.1,7.199999999999999) {边界：$AF=8/3$，已占满总量};
\draw[gray,dashed] (2.2,5.85) rectangle (8.2,6.75);\fill[blue!30] (2.2,6.05) rectangle (8.2,6.55);
\draw[violet,dashed] (8.2,5.8)--(8.2,6.8999999999999995);\node[anchor=west] at (10.7,6.3) {余量为零};
\node at (5.2,5.5) {$1/BF=0$：没有有限另一段};
\node at (7.1,3.65) {例4：假设同支 $AF=2$，已超过总量};
\draw[gray,dashed] (2.2,2.3) rectangle (8.2,3.2);\fill[blue!30] (2.2,2.5) rectangle (10.2,3.0);
\draw[violet,dashed] (8.2,2.25)--(8.2,3.35);\node[anchor=west] at (10.7,2.75) {不可能有正余量};
\node at (5.2,1.95) {$1/2>3/8$：所设同支情形无解};
\node at (7.1,.55) {虚框宽度表示 $K=3/8$；色条画的是倒数，不是原线段长度。};
\end{tikzpicture}
\end{center}
  \textbf{\textcolor{CExample}{例4（迁移：公式求出的“负长度”是在提醒无解）}}

  \textbf{题目：} 在 $\dfrac{x^2}{9}-\dfrac{y^2}{16}=1$ 中，能否找到\textbf{同在右支}的两个焦点直线交点 $A,B$，使 $|AF|=2$？

  \textbf{解：} 同支公式要求
  \[
    \frac12+\frac1{|BF|}=\frac38
    \quad\Longrightarrow\quad
    \frac1{|BF|}=-\frac18<0,
  \]
  与 $|BF|>0$ 矛盾，故\textbf{不存在}。虽然右顶点 $(3,0)$ 到 $F(5,0)$ 的距离正好是 $2$，但它所在的水平焦点直线另一有限交点在\textbf{左支}，而不是右支。
  \[
    \boxed{\text{所述同支情形不存在}}.
  \]
  \textbf{记忆检查：} 同支时每段必须满足 $r>p/2=b^2/(2a)$（由另一个倒数为正得到）。

  \medskip
\end{examplebox}\clearpage\begin{examplebox}
\begin{center}\begin{tikzpicture}[font=\small,>=stealth]
\draw[rounded corners=3pt,gray!45] (0,0) rectangle (14.2,12);
\node at (7.1,11.45) {先数有限交点：斜率不存在，与仅有一个交点是两回事};\node at (2.35,10.2) {水平：异支两点};\begin{scope}[shift={(2.35,6)},scale=.18]
\begin{scope}\clip (-10,-9) rectangle (10,9);
\draw[teal!75!black,thick,samples=110,domain=-1.9:1.9,variable=\t] plot ({3*cosh(\t)},{4*sinh(\t)});
\draw[teal!75!black,thick,samples=110,domain=-1.9:1.9,variable=\t] plot ({-3*cosh(\t)},{4*sinh(\t)});
\draw[blue!75!black,very thick] (-10,0)--(10,0);\end{scope}\draw[->,gray] (-10.5,0)--(10.5,0);\fill (5,0) circle(8pt);\node at (1,-1.5) {$F$};\draw[gray] (2.5,-1.0)--(4.8,-.05);\fill[red!75!black] (-3.0000000000,0.0000000000) circle(10pt);\fill[red!75!black] (3.0000000000,0.0000000000) circle(10pt);
\end{scope}
\node[align=center] at (2.35,2.7) {$k=0$\\[13pt]可以用倒数差};
\node at (7.1,10.2) {平行渐近线：仅一点};\begin{scope}[shift={(7.1,6)},scale=.18]
\begin{scope}\clip (-10,-9) rectangle (10,9);
\draw[teal!75!black,thick,samples=110,domain=-1.9:1.9,variable=\t] plot ({3*cosh(\t)},{4*sinh(\t)});
\draw[teal!75!black,thick,samples=110,domain=-1.9:1.9,variable=\t] plot ({-3*cosh(\t)},{4*sinh(\t)});
\draw[gray,dashed] (-8,-10.6667)--(8,10.6667);\draw[blue!75!black,very thick] (-2,-9.3333333)--(12,9.3333333);\end{scope}\draw[->,gray] (-10.5,0)--(10.5,0);\fill (5,0) circle(8pt);\node at (1,-1.5) {$F$};\draw[gray] (2.5,-1.0)--(4.8,-.05);\fill[red!75!black] (3.4000000000,-2.1333333333) circle(10pt);
\end{scope}
\node[align=center] at (7.1,2.7) {$k=b/a=4/3$\\[13pt]不能虚设第二点};
\node at (11.85,10.2) {竖直：同支两点};\begin{scope}[shift={(11.85,6)},scale=.18]
\begin{scope}\clip (-10,-9) rectangle (10,9);
\draw[teal!75!black,thick,samples=110,domain=-1.9:1.9,variable=\t] plot ({3*cosh(\t)},{4*sinh(\t)});
\draw[teal!75!black,thick,samples=110,domain=-1.9:1.9,variable=\t] plot ({-3*cosh(\t)},{4*sinh(\t)});
\draw[blue!75!black,very thick] (5,-9)--(5,9);\end{scope}\draw[->,gray] (-10.5,0)--(10.5,0);\fill (5,0) circle(8pt);\node at (1,-1.5) {$F$};\draw[gray] (2.5,-1.0)--(4.8,-.05);\fill[red!75!black] (5.0000000000,5.3333333333) circle(10pt);\fill[red!75!black] (5.0000000000,-5.3333333333) circle(10pt);
\end{scope}
\node[align=center] at (11.85,2.7) {$x=c=5$\\[13pt]可以用倒数和};
\node at (7.1,1.1) {三图同一双曲线、同尺度；红点为有限交点，灰虚线为渐近线。};\node at (7.1,.45) {$k=-b/a$ 与中图上下对称，也只有一个有限交点。};
\end{tikzpicture}
\end{center}
  \textbf{\textcolor{CExample}{例5（边界：渐近线斜率不是双交点）}}

  \textbf{题目：} 对双曲线 $\dfrac{x^2}{9}-\dfrac{y^2}{16}=1$，过右焦点 $F(5,0)$ 的直线 $y=\dfrac43(x-5)$，能否直接应用焦点分弦定理？

  \textbf{解：} 这里 $k=4/3=b/a$，联立后的 $x^2$ 系数 $b^2-a^2k^2=0$，只有一个有限交点
  \[
    P\left(\frac{17}{5},-\frac{32}{15}\right).
  \]
  因而\textbf{不满足“两个不同的有限交点”}，不能虚设另一交点，也不能机械套用加减公式。
  \[
    \boxed{\text{不适用；此直线仅有一个有限交点。}}
  \]
\end{examplebox}

\clearpage
\begin{examplebox}
  \textbf{\textcolor{CExample}{例6（焦准配对：A 在左支，仍然可以选右焦点与右准线）}}

  \textbf{题目：} 双曲线 $\dfrac{x^2}{9}-\dfrac{y^2}{16}=1$，
  左顶点 $A(-3,0)$、右顶点 $B(3,0)$、右焦点 $F_2(5,0)$。
  直线 $AB$ 经过右焦点 $F_2$（注意 $F_2$ 在\textbf{线段延长线}上）。
  对点 $A,B$，分别检验\textbf{右焦点--右准线}的距离比，并顺便检验 C057 异支公式。

  \textbf{第一步（固定一套配对）：} $a=3,b=4,c=5,e=5/3$；
  右焦点 $F_2(5,0)$ 对应\textbf{右准线} $\ell_2:x=a^2/c=9/5$。

  \textbf{第二步（左支点 $A$，仍用右准线）：}
  \[
    |AF_2|=|-3-5|=8,\qquad
    d(A,\ell_2)=\left|-3-\frac95\right|=\frac{24}{5}.
  \]
  \[
    \boxed{\frac{|AF_2|}{d(A,\ell_2)}=\frac{8}{24/5}=\frac53=e}.
  \]

  \textbf{第三步（右支点 $B$，同一套配对）：}
  \[
    |BF_2|=|3-5|=2,\qquad
    d(B,\ell_2)=\left|3-\frac95\right|=\frac65.
  \]
  \[
    \boxed{\frac{|BF_2|}{d(B,\ell_2)}=\frac{2}{6/5}=\frac53=e}.
  \]
  \textbf{再看 C057：} $A,B$ 分居左右两支，
  \[
     \left|\frac1{|AF_2|}-\frac1{|BF_2|}\right|
     =\left|\frac18-\frac12\right|=\frac38
     =\frac{2a}{b^2}.
  \]
  \textbf{结论：}左支点 $A$ \textbf{无需换左焦点或左准线}；
  必须保证\textbf{焦点与准线本身配对}。例如拿 $AF_2$ 去除以
  $A$ 到\textbf{左准线} $x=-9/5$ 的距离，就不是 $e$：
  \[
    \frac{|AF_2|}{d(A,\ell_1)}=\frac8{6/5}=\frac{20}{3}\ne\frac53.
  \]
  \textbf{一句话记忆：焦准成对，取点自由；比值是 $e$，不是 $1$。}
\end{examplebox}


\clearpage
\begin{examplebox}
  \textbf{\textcolor{CExample}{例7（消除疑惑：水平线 $\alpha=0$，为什么右支公式算出 $-8$？）}}

  \textbf{题目：} 对 $\displaystyle\frac{x^2}{9}-\frac{y^2}{16}=1$，
  $F=(5,0)$。把水平直线的倾斜角 $\alpha=0$ 代入
  $\displaystyle r=\frac{b^2}{a-c\cos\varphi}$ 得到 $-8$，
  这是否意味着存在负的焦半径？

  \textbf{解法一（最低门槛：先看点的位置）。}
  双曲线的两个顶点是 $(3,0)$、$(-3,0)$，均在 $F=(5,0)$ 左方。
  \textbf{直线倾斜角}为 $\alpha=0$，但从 $F$ 指向两顶点的
  \textbf{射线方向角}均为 $\varphi=\pi$，所以 $\cos\varphi=-1$。
  \[
    \begin{aligned}
      \text{右支：}\quad r_{\rm R}&=\frac{16}{3-5(-1)}=2,\\[3pt]
      \text{左支：}\quad r_{\rm L}&=-\frac{16}{3+5(-1)}=8.
    \end{aligned}
  \]
  若把 $\alpha=0$ 冒充 $\varphi=0$，右支式算出的 $16/(3-5)=-8$
  \textbf{不满足}右支射线公式 $a-c\cos\varphi>0$ 的适用条件。

  \textbf{解法二（快速验证：有向位移）。}
  取整条直线的正向为 $(1,0)$，设 $(x,y)=(5+t,0)$。则
  \[
    \frac{(5+t)^2}{9}=1
    \ \Longrightarrow\ (t+2)(t+8)=0
    \ \Longrightarrow\ \boxed{t=-2,\;-8}.
  \]
  这里 $t=-2$ 对应右顶点 $(3,0)$，$t=-8$ 对应左顶点 $(-3,0)$。
  \textbf{负的是有向位移，不是距离。}

  \textbf{联结 C057：}两交点异支，正长度为 $2$ 与 $8$，因此
  \[
    \boxed{\frac12-\frac18=\frac38=\frac{2a}{b^2}}.
  \]
  \textbf{作答检查：}先辨\textbf{直线角/射线角}，再辨\textbf{分支}，
  最后才将有向参数改写为正距离。
\end{examplebox}
